% Compilação: `make` dentro de writing/, ou manualmente:
%   latexmk -pdf -shell-escape -interaction=nonstopmode -synctex=1 \
%           -auxdir=.aux_files -bibtex main.tex
%
% Classe: abntex2 (cópia local em setup/, não empacotada no TeX Live desta
% máquina — kpsewhich abntex2.cls falhou; verificado em 2026-09-08).
% Citações: biblatex + estilo abnt (pacote biblatex-abnt, backend biber),
% já presente no TeX Live do sistema. abntex2cite NÃO está instalado e não
% foi usado.
%
% Orçamento de página do relatório inteiro (4-7 páginas incl. referências,
% ver enunciado Seção 11): ver comentários de orçamento em cada seção.

\documentclass[article,11pt,oneside,a4paper,brazil,section=TITLE]{setup/abntex2}

\usepackage[T1]{fontenc}
\usepackage[utf8]{inputenc}
\usepackage{lmodern}
\usepackage{microtype}
\usepackage{amsmath}
\usepackage{amssymb}
\usepackage{graphicx}
\usepackage{booktabs}
\usepackage{tabularx}
\usepackage[backend=biber,style=abnt]{biblatex}
\usepackage{csquotes}
\addbibresource{referencias.bib}
\nocite{wollschlager2025geometry}

\graphicspath{{pictures/}}
\hypersetup{hidelinks}
\usepackage{xurl}

% Margens levemente reduzidas (ABNT recomenda 3-2-2-3) para caber o
% conteúdo em 4-7 páginas incl. referências, exigido pelo enunciado.
\setlrmarginsandblock{2cm}{2cm}{*}
\setulmarginsandblock{2cm}{2cm}{*}
\checkandfixthelayout
\setlength{\bibitemsep}{0.3\baselineskip}

\newcolumntype{Y}{>{\raggedright\arraybackslash}X}

\titulo{Análise Crítica: A Geometria da Recusa em Modelos de Linguagem de Grande Escala}
\autor{Gabriel Salmoria}

\begin{document}
\textual

\begin{center}
  {\small
    UNIVERSIDADE FEDERAL DE SANTA CATARINA\\
    Departamento de Informática e Estatística\\
    INE5448 --- Inteligência Artificial e Segurança da Informação\par}
  \vspace{1.2em}
  {\bfseries\large Análise Crítica de Artigo Científico:\\
   A Geometria da Recusa em Modelos de Linguagem de Grande Escala\par}
  \vspace{0.8em}
  {Gabriel Salmoria \quad --- \quad \today\par}
\end{center}
\vspace{1em}

% =====================================================================
% Orçamento: seção curta (poucas linhas).
\section{Identificação do artigo}
\label{sec:identificacao}

\begin{description}
  \item[Autores] Tom Wollschläger, Jannes Elstner, Simon Geisler, Vincent
    Cohen-Addad, Stephan Günnemann, Johannes Gasteiger.
  \item[Título] The Geometry of Refusal in Large Language Models: Concept
    Cones and Representational Independence.
  \item[Veículo] \textsc{Proceedings of the 42nd International Conference
    on Machine Learning} (ICML 2025), Vancouver, Canadá.
    \textsc{Proceedings of Machine Learning Research}, v.~267,
    p.~66945--66970. Disponível em:
    \url{https://proceedings.mlr.press/v267/wollschlager25a.html}.
  \item[Ano] 2025.
  \item[DOI] Nenhum. A PMLR não atribui DOI aos artigos publicados em suas
    atas (landing page conferida em 2026-09-08).
  \item[Classificação do veículo] ICML tem classificação CORE A* (\url{https://portal.core.edu.au/conf-ranks/1121/}, ranking ICORE2026, conferido em 2026-09-08). Qualis (CAPES/Sucupira): NÃO VERIFICADO, não consultei a Plataforma Sucupira diretamente.
\end{description}

Existe também uma versão no arXiv (2502.17420v2), posterior ao camera-ready; uma comparação entre os dois PDFs não achou apêndices extras nela, só uma diferença pontual, já que a v2 rotula como \emph{Theorem 4.1} algo que na PMLR é \emph{Definition 4.1} (as propriedades de recusa do RCO, Seção~\ref{sec:proposta}).

% =====================================================================
% Orçamento: seção curta.
\section{Justificativa da escolha}

O tema deste artigo, mecanismos de recusa e jailbreaking em LLMs, atende aos critérios de seleção da atividade e se conecta bem ao eixo de IA e segurança da disciplina. O veículo tem classificação A* (Seção~\ref{sec:identificacao}) e é um artigo de 2025, o que faz ele se mostrar atual e bem colocado na discussão de explicabilidade de modelos de LLM. O tema discutido sendo recusa de modelos, tecnicas de jailbreaking e sandboxing tem se tornado cada vez mais proeminentes na discussão academica, uma vez que modelos de IA, como o GPT-5.6 da openAI tomou atitudes como a invasão de servidores e ambientes da HuggingFace durante testes de benchmark, ainda nesse ano.

% =====================================================================
% Orçamento (c-g juntas): ~2,5 páginas no total.
\section{Contextualização e fundamentação}

A ideia de que coceitos abstratos, como sentimento, gênero, tamanho, e, nesse contexto, recusa, correspondem a direções lineares no espaço de ativação de uma LLM vem de uma linha de pesquisa apoiada na hipótese da representação linear \parencite{park2023linear} e sistematizada como ferramenta de análise por Representation Engineering \parencite{zou2023representation}, de onde o artigo herda as operações de ablation e adição de ativação (Seção~\ref{sec:proposta}). O antecedente mais direto é \textcite{arditi2024refusal}, que usa esse tipo de direção para defender que a recusa é mediada por uma única direção linear, resultado influente citado inclusive por trabalhos que tomam essa direção como alvo de ataques adversariais \parencite{yu2024robust}. É também o que este artigo contesta.

Um outro trabalho, \textcite{pan2025hidden} (ICML 2025, não preprint), chega a uma conclusão parecida por um caminho diferente, analisando como a representação de segurança muda durante o ajuste fino de alinhamento. Os próprios autores citam esse trabalho como evidência independente a favor da sua tese (retomado na Seção~\ref{sec:comparacao}).

% =====================================================================
\section{Problema de pesquisa}

O artigo pergunta, em três formas diferents, se um método baseado em gradiente encontra direções de recusa melhores que a diferença de médias usada até então, se essa recusa vem mesmo de uma única direção linear ou de algo mais complexo/elaborado no espaço semântico da rede neural em si, e se existem direções independentes entre si, acionáveis diretamente pelo texto de entrada, sem mexer nas ativações internas.

A premissa em si desafiada é a de \textcite{arditi2024refusal}, de que uma única direção linear é responsável pela recusa, e o escopo fica restrito a modelos abertos, com acesso \emph{white-box} às ativações e porte moderado, até 14 bilhões de parâmetros (Seção~\ref{sec:metodologia}). É claro, ficam de fora os modelos fechados e proprietarios, nos quais esse tipo de analise não é possível.

% =====================================================================
\section{Proposta dos autores}
\label{sec:proposta}

Métodos anteriores e mais comuns na área, como o de \textcite{arditi2024refusal}, calculam a direção de recusa pela diferença entre a ativação média do modelo em prompts nocivos e inofensivos, a direção \emph{difference-in-means} (DIM), sem nenhuma otimização voltada ao efeito que deve produzir na prática.

O RDO (\emph{Refusal Direction Optimization}) ajusta essa direção $r$ por gradiente em vez de estimá-la, até que, removida das ativações, o modelo volte a responder pedidos antes recusados e, somada a prompts inofensivos, passe a recusá-los, com um terceiro objetivo preservando o resto do comportamento; os três entram no treino como termos de perda otimizados ao mesmo tempo.

Em vez de assumir que existe uma única direção de recusa, o passo seguinte, RCO, treina um conjunto de vetores ortogonais que juntos formam uma região do espaço de ativação, um cone poliédrico onde qualquer direção satisfaz as mesmas propriedades do RDO.

Mesmo treinando uma direção RDO ortogonal à DIM, os autores notam que remover uma do modelo ainda afeta o quanto a outra aparece nas camadas seguintes, a observação inesperada por trás da independência representacional (RepInd). Como ortogonalidade não garante independência de causalidade nos modelos, o critério que propõem é em si mais rigoroso, sendo duas direções que só contam como independentes se remover uma não mudar o quanto a outra permanece presente nas ativações, camada a camada.

% =====================================================================
\section{Metodologia experimental}
\label{sec:metodologia}
% Orçamento total de c-g: ~2,5 páginas -- ajustar tamanho de e+f se
% c, d, g (acima/abaixo) crescerem além do esperado quando preenchidas.

A avaliação usa sete modelos de chat de até 14 bilhões de parâmetros, das famílias Qwen2.5, Gemma 2 e Llama 3, todos com acesso \emph{white-box} às ativações, e os dados vêm de bancos padrão de prompts nocivos e inofensivos (SALAD-Bench e Alpaca). Um detalhe que chama atenção no protocolo é que os alvos usados pra treinar o RDO vêm da própria direção DIM, não de um modelo externo independente.

A Tabela~\ref{tab:hparams} resume os principais hiperparâmetros do treino.

\begin{table}[htb]
  \centering\small
  \caption{Hiperparâmetros principais reportados no artigo.}
  \label{tab:hparams}
  \begin{tabularx}{\textwidth}{@{}lY@{}}
    \toprule
    \textbf{Parâmetro} & \textbf{Valor} \\
    \midrule
    Otimizador & AdamW, weight decay 0 \\
    Taxa de aprendizado base & 0,01; reduzida em plateau a cada 5 batches (÷10, até 2x) \\
    Batch size total / grad. accumulation & 16 / 16 passos \\
    $\lambda_{\text{abl}}$, $\lambda_{\text{add}}$, $\lambda_{\text{ret}}$ (RDO) & 1,0 / 0,2 / 1,0 \\
    Amostras Monte Carlo (RCO) & 16 por passo de acumulação (256 efetivas/batch) \\
    $\lambda_{\text{ind}}$, corte de camadas (RepInd) & 200; primeiros 90\% das camadas \\
    \bottomrule
  \end{tabularx}
\end{table}

O sucesso de jailbreak é medido pela taxa de ataques bem-sucedidos (ASR) em vários benchmarks padrão da área, com um juiz automático, e efeitos colaterais são medidos em benchmarks gerais de capacidade do modelo, como TruthfulQA. Código e hiperparâmetros estão publicados, mas nenhuma seção reporta custo computacional, o que é uma pena porque impede comparar o custo de treinar uma direção por RDO com o de calcular um DIM, que é praticamente de graça.

% =====================================================================
% Orçamento: ~0,75 página.
\section{Resultados e discussão}
\label{sec:resultados}

Na prática, sob o tipo de intervenção mais direto, que remove a direção das ativações, RDO e DIM acabam tendo desempenho parecido nos benchmarks de jailbreak; a vantagem do RDO aparece mais em efeitos colaterais menores sobre capacidades gerais do modelo, com uma redução considerável do erro induzido em benchmarks como o TruthfulQA, do que na taxa de sucesso de ataque em si. A extensão para múltiplas direções (RCO) funciona até certo ponto: cones de mais dimensões existem em todos os modelos testados, mas modelos maiores sustentam cones maiores, relação que os próprios autores tratam como hipótese plausível, não como algo derivado formalmente.

Os experimentos voltados à independência representacional (RepInd) apontam na mesma direção geral, incluindo evidência de que essas direções são acessíveis a partir de manipulações no próprio texto de entrada, não só via intervenção direta nas ativações. Nem todos esses testes, porém, foram repetidos no conjunto inteiro de sete modelos, e nenhuma das tabelas de resultado traz intervalo de confiança ou múltiplas execuções, apesar de os próprios autores relatarem variância alta ao treinar essas direções (Seção~\ref{sec:metodologia}).

% =====================================================================
% Orçamento: ~0,75 página.
\section{Comparação com a literatura}
\label{sec:comparacao}

O contraste mais direto é mesmo com \textcite{arditi2024refusal} (NeurIPS 2024), já que a conclusão central deles, de que uma única direção linear medeia toda a recusa, é exatamente o que esse artigo revisita. A diferença entre os dois é mais conceitual do que numérica (Seção~\ref{sec:resultados}), mas é o que permite estender o método pra múltiplas direções (RCO) e testar se elas são independentes (RepInd), coisa que Arditi et al. não têm.

\textcite{pan2025hidden} (ICML 2025) chega a uma conclusão parecida, múltiplas direções em vez de uma só, mas por um caminho diferente, analisando como a representação de segurança muda durante o ajuste fino, em vez de propor um método de descoberta para qualquer modelo já treinado. O repertório de intervenções em si vem de \textcite{zou2023representation} e é usado sem alteração por DIM, RDO, RCO e RepInd; a contribuição deste artigo está em como a direção é obtida e validada, não nas intervenções.

Fora das referências do próprio artigo, \textcite{yang2026beyond} (ainda preprint) testa a mesma pergunta em modelos de raciocínio com chain-of-thought, e encontra que a ativação sozinha reverte a recusa em só 39\% dos casos quando a cadeia de raciocínio é mantida, contra 70\% quando é removida, sugerindo um limite para toda essa linha, incluindo RDO/RCO/RepInd: em modelos com raciocínio explícito, a recusa também pode não estar só no espaço de ativação residual.

% =====================================================================
% Orçamento: curta.
\section{Pontos fortes}

O que mais impressiona é a cobertura experimental, com sete modelos de três famílias, três benchmarks de jailbreak e quatro de capacidade geral, reduzindo bastante a chance de o resultado depender de um único modelo ou métrica -- essa é a base concreta por trás do principal mérito do artigo, que é desafiar, com um método sistemático, uma conclusão citada com frequência na literatura de interpretabilidade e segurança de LLMs. Hiperparâmetros completos, código público e uma quantidade grande de figuras e tabelas nos apêndices tornam boa parte das afirmações verificável direto no artigo.

% =====================================================================
% Orçamento: ~1,5 página.
\section{Limitações e críticas}

O método dos autores usa a própria direção de base (DIM) para gerar os alvos que treinam o RDO (Seção~\ref{sec:metodologia}), o que é diferente da limitação que eles já admitem (Seção 7 do artigo) sobre os alvos serem iguais entre os métodos que propõem. O problema aqui é que a DIM, usada como referência na comparação, também ajuda a supervisionar o treino do RDO, e por isso o empate entre os dois na prática (Seção~\ref{sec:resultados}) pode não significar que acharam uma direção melhor -- só que o RDO aprendeu a imitar a DIM, já que os autores não testam alvos gerados de outra forma para descartar essa possibilidade.

As tabelas de resultado trazem só um número por combinação de modelo e método, sem intervalo de confiança nem repetição com sementes diferentes, e isso pesa mais porque os próprios autores relatam variância alta ao treinar as direções mais restritivas, as RepInd (Seção~\ref{sec:metodologia}); se uma única direção já varia tanto, as diferenças de um ou dois pontos entre os dois métodos na Tabela 5 podem ser só ruído de uma execução. O teto de escala também incomoda, já que o artigo vai até 14 bilhões de parâmetros, contra os 13 modelos de até 72 bilhões que Arditi et al. (2024) testaram.

Exigir que uma medida geométrica, o ângulo entre a direção e a ativação, não mude quando a outra é removida não deixa claro se isso basta para garantir independência de verdade, no sentido de causa e efeito -- e essa é justamente a forma como os autores definem independência (Seção~\ref{sec:proposta}). O experimento mais próximo de testar isso roda só no Gemma~2~2B e não compara com um controle óbvio, como direções aleatórias quaisquer, de modo que o ganho observado pode vir simplesmente de remover mais coisa do modelo, não da independência em si.

Na prática, a remoção sempre atua sobre uma única direção por vez, mesmo quando ela vem de um cone com várias dimensões possíveis -- um detalhe técnico do cone de direções (Seção~\ref{sec:proposta}) que os autores nunca testam contornar removendo o espaço todo de uma vez só. Fica a dúvida se ter várias dimensões muda alguma coisa de fato, ou se só dá mais opções de direção única pra escolher.

O procedimento mais caro do artigo, o de treinar as direções RepInd, exige 25 candidatos por direção, e nenhuma seção reporta quanto isso custa em termos computacionais (Seção~\ref{sec:metodologia}); sem esse número não dá pra comparar com o custo de calcular a DIM, que é praticamente de graça.

Por fim, o Impact Statement do artigo, que diz que o foco em acesso \emph{white-box} torna os achados ``principalmente aplicáveis à defesa'', entra em tensão com a própria Seção 6.2, que usa esse mesmo acesso para treinar sufixos de texto que jailbreakam o modelo em 36\% das tentativas. Um sufixo, uma vez encontrado, não precisa mais desse acesso pra ser usado, só pra ser descoberto, e o artigo não testa se ele funciona em outros modelos.

% =====================================================================
% Orçamento: ~0,5 página.
\section{Trabalhos futuros}

Duas extensões dão pra testar pontos da seção anterior. Uma ataca a circularidade dos alvos: se o empate entre RDO e DIM vem de os alvos serem gerados pela própria direção DIM, treinar o RDO com alvos independentes (respostas humanas, ou geradas por outro modelo) deve abrir uma diferença maior de ASR -- o experimento mínimo é repetir o treino em um ou dois modelos trocando só a fonte dos alvos, mantendo o resto do protocolo (Seção~\ref{sec:metodologia}).

A outra testa se o cone do RCO faz algo além de ampliar o menu de direções de posto~1: remover o subespaço inteiro gerado pela base $B$ deveria produzir ASR diferente do observado ao remover a melhor direção única do cone, o que dá para comparar nos mesmos modelos e benchmark usados na Figura~3 do artigo.

% =====================================================================
% Orçamento: curta.
\section{Uso de ferramentas de Inteligência Artificial}

Usamos o Claude para montar o esqueleto LaTeX do relatório, e redigir um primeiro rascunho de todas as seções. Claro, esse rascunho foi apenas uma base, e passou por revisão minha antes da entrega, com uma reescrita geral do texto. (Uso assim, apenas pra ter uma base pra construir o texto, quais ideias expor e etc, uso de IA pra fazer um "esqueleto" mais bem montado)

% =====================================================================
% Orçamento: curta.
\section{Conclusão}

Esse trabalho me deixou bem mais desconfiado de alegações de que um comportamento complexo de um LLM é controlado por uma única direção no espaço de ativação, e reforçou a importância de separar resultado reportado de explicação post-hoc dos autores. A parte mais relevante para minha formação em IA e Segurança foi perceber como decisões de protocolo aparentemente neutras, como reaproveitar a direção DIM para gerar os próprios alvos de treino do método que compete com ela, podem comprometer a força de uma conclusão.

% =====================================================================
\printbibliography[title=Referências]

\end{document}
